\chapter*{Preface}
\addcontentsline{toc}{chapter}{Preface}
% Front matter, on the operator's directives (2026-07-25 15:14 + 15:18):
% the case for MANY being unnecessary -- the redundancy as a consequence of
% representation, not of the world appearing 3+1 dimensional. Laws: the
% many-earn per Kodo's bar (earned like Vol 2's 4, never asserted);
% least-squares cited with the Legendre/Gauss attribution told honestly;
% structural-vs-statistical held; the anti-continuum tie; the
% representation-not-world fence IN THE SAME BREATH (Beastmaster verifies
% on arrival, founding-program-block protocol); the INTUITION LAW stated
% with Beastmaster's adds-no-claim tightening. Gate: Kodo (hardest on the
% earn) + Beastmaster (the fence).
% EDIT 2026-07-26 (operator "preface 3 needs to define 5 as the thing that
% gives no information," confirmed "yes (not the target)"): the FIVE
% crystallized at the end of the overdetermination paragraph — "five is the
% first reading that gives no information," scoped in the same breath
% (confirm-not-inform / nothing remains to eliminate), earned off the page's
% OWN least-squares/overdetermination argument. NO device machinery (no
% T=5 / de-selection / mediant — operator confirmed not-the-target); inside
% the representation-not-world fence; adds-no-claim (names what the page
% already argued). Gate: Kodo (earn) + Beastmaster (arrival-read hardest
% against device-import + world-fence trip) → re-publish.

The last volume's preface earned the number four: three soundings leave a
point and its mirror twin, and the fourth, off-plane, breaks the tie. This
preface asks the question a careful reader asked next, and answers it,
because the answer is this volume's first lesson in how a representation
behaves. Why stop at four? Why not five soundings, or six, or a hundred?
Instruments are cheap; take \emph{many}. The answer --- and it deserves to
be earned, not waved --- is that many buys nothing, and the reason it buys
nothing is a fact about representation itself.

Watch the fifth sounding arrive. Four have done their work: the candidates
were cut to two, the tie was broken, the point has been named. Now the
fifth reading comes in --- and there is nothing left for it to say. If the
apparatus is honest, the fifth reading must pass through the point the
four already fixed; its value is implied before it is taken. It can
\emph{confirm} --- and confirmation has real uses --- but it cannot
\emph{inform}: no candidate remains for it to eliminate. And what holds
for the fifth holds for the fiftieth. Once a representation with three
coordinates and one tiebreak has been filled, every further constraint is
spoken for in advance. The trade's word for this is overdetermination, and
its classical mathematics is not folklore but a named subject: the method
of least squares, published by Legendre in 1805 and given its foundation
by Gauss --- who had used it privately before the publication, and who
built its theory in 1809 and again in 1821 --- precisely for the
overdetermined systems of astronomy and geodesy, where observations
outnumber unknowns as a matter of course. The reader should note what the
method \emph{does} with the surplus: it does not extract new coordinates
from it. It manages error. The honest job of the pile has been, since the
year it was named, consistency and noise --- never information. And since
this page has now earned it, the number may as well receive its
definition here, where it stands: \emph{five is the first reading that
gives no information} --- the first constraint spoken for in advance,
able to confirm what the filled representation has already named and
unable to eliminate anything, because nothing remains to eliminate. Four
is the count that completes the representation; five is the name of the
surplus.

That distinction --- noise against structure --- is the earn's load-bearing
beam, and the last volume already set it. Many measurements are the
correct response to \emph{noise}: readings scattered by accident average
toward truth, and the pile genuinely helps. But the two-candidate
ambiguity of trilateration was never noise; the last preface said so in so
many words --- it is structural, it lives in the geometry, not the noise
--- and against structure the pile is helpless. A thousand repetitions of
the same three soundings polish the same two candidates to a brilliant
shine and choose neither. What broke the tie was not more of the same; it
was one measurement of a \emph{different character}, placed off the plane.
Structure yields to structure, never to volume. So the case for many
collapses into two honest halves: against noise, many works and is
welcome; against structure, many is spoken for in advance --- and the
naming of a point in a filled representation is structure through and
through. The representation was complete at four, and completeness is
indifferent to enthusiasm.

The reader who has followed this series will hear a familiar chord. The
drift toward many --- more soundings, more repetitions, more and finer
forever --- is the drift toward the continuum, and this series has held a
fence against that drift since its first volume: the counts that matter
have been small and structural throughout, and the pile has never once
been where the content lived. This preface adds the surveyor's case to
that standing discipline and claims nothing further from it.

\begin{fence}
And now the fence, in the same breath, because it is the load-bearing half
of everything above. Nothing in this preface is a claim about the world
--- not about its dimensions, not about what space is, not about how many
directions reality possesses. The four that sufficed and the many that is
unnecessary are facts about a \emph{representation}: about how many
independent constraints fill a description that carries three coordinates
and one tiebreak. A reader could build a representation with more
coordinates or fewer and the arithmetic would move with the
representation, which is exactly the point --- the count belongs to the
description, and the world is left un-claimed. This series said on the
first page of its first volume that the electron, in these books, is the
name of a mathematical model; the same licence governs here, applied to
dimensionality. Whoever reads the four of the last preface, or the
completeness of this one, as this series measuring the dimensions of
spacetime has taken a step no volume of this series takes.
\end{fence}

One more law belongs on this first page, because this volume is licensed
to do something its predecessor forbade, and a licence without a law is a
leak. The last volume permitted only cited ideas; this one is called
\emph{The Intuition}, and it will hand the reader pictures --- the soap
film, the mountain pass, the bookkeeper's zero, the two hands closing ---
because pictures are how proof structures are \emph{felt} before they are
checked. The law is this: every intuition in this book is an aid, and it
is always marked as one --- \emph{intuition offered, not claim made} ---
and no intuition may add a single mathematical claim that its cited
sources do not already contain. An intuition here is a way of seeing cited
material, never new content; the back matter will keep a ledger of every
intuition beside the citation it re-sees; and the failure this book hunts
in itself is the intuition that quietly hardens into an assertion. Claims
keep the last volume's law --- a name and a date, or they do not appear.
Pictures keep this one's --- a label, or they do not appear.

With the number earned, the fence set, and the law posted, the volume can
say in one sentence what it is for: four faces of mathematics ---
geometry, topology, algebra, and numerics --- each read as a structure a
variational proof can take, so that the constructions of the last volume
and the measurements of the first stop being merely verified and start
being obvious. Read on with the label in mind.
